\begin{titlepage}
\thispagestyle{empty}
% Magazine "image-hero" cover: the real author-collaboration network is a large
% offset hero image; the title layers over it; metadata is right-aligned to break
% the single left axis. All placement is on an absolute mm grid for art direction.
\begin{tikzpicture}[remember picture,overlay,x=1mm,y=1mm,shift={(current page.north west)}]
  \fill[white] (0,0) rectangle (210,-297);
  % Large author network as the hero image — near full-bleed, anchored to the
  % lower-right so it dominates the middle-lower page and the stat layers over it.
  \node[opacity=0.20,anchor=north east,inner sep=0] at (218,-86)
    {\includegraphics[width=214mm]{author_collaboration_network_cover.png}};

  % Masthead: ink wordmark (left), right-aligned issue meta, gold hairline.
  \node[anchor=west,inner sep=0,font=\coverSans\bfseries\LARGE] at (17,-14)
    {\textcolor{SciGold}{$\blacklozenge$}\,\textcolor{SciInk}{SciScope}};
  \node[anchor=east,inner sep=0,font=\coverCJK\footnotesize,text=SciMuted] at (193,-13.4)
    {数据分析报告 · 交付物 · 公开科技文献 2022–2026};
  \draw[SciGold,line width=0.9pt] (17,-20) -- (193,-20);

  % Title block (left), gold kicker stroke, layered over the network.
  \draw[SciGold,line width=2.6pt] (17.5,-72) -- (33,-72);
  \node[anchor=north west,inner sep=0,font=\coverTitle\fontsize{62}{68}\selectfont,text=SciInk] at (16,-80)
    {数据分析报告};
  % Deck — RIGHT aligned, two lines, to set up asymmetry against the title.
  \node[anchor=north east,inner sep=0,align=right,font=\coverCJK\fontsize{15}{21}\selectfont,text=SciTealDark] at (193,-110)
    {近五年公开科技文献的\\分布、主题演化与合作网络};

  % Hero stat (left, Bodoni) with gold underline; caption RIGHT aligned alongside.
  \node[anchor=north west,inner sep=0,font=\coverSans\footnotesize,text=SciMuted] at (17,-150)
    {CORPUS AT A GLANCE \;/\; 核心规模};
  \node[anchor=north west,inner sep=0,font=\coverNum\bfseries\fontsize{37}{39}\selectfont,text=SciTealDark] at (16,-156)
    {168,447};
  \draw[SciGold,line width=1.4pt] (17,-172) -- (45,-172);

  % Coverlines — the report's facets, like a magazine's cover lines, set low and
  % overlaid on the hero image.
  \node[anchor=west,inner sep=0,font=\coverSans\scriptsize,text=SciGold] at (17,-236)
    {IN THIS REPORT};
  \node[anchor=north west,inner sep=0,font=\coverCJK\large,text=SciTealDark] at (17,-243)
    {文献分布\goldsep 关键词演化\goldsep 作者合作网络\goldsep 数据质量边界};

  % Footer: hairline + muted scope notes (left object / right boundary).
  \draw[SciLine,line width=0.5pt] (17,-280) -- (193,-280);
  \node[anchor=west,inner sep=0,font=\coverCJK\scriptsize,text=SciMuted] at (17,-286)
    {分析对象:公开科技文献元数据 · 摘要 · 作者 · 关键词 · 年份 · 来源};
  \node[anchor=east,inner sep=0,font=\coverCJK\scriptsize,text=SciMuted] at (193,-286)
    {结论限定在公开数据 · 近五年窗口 · 当前字段完整度之内};
\end{tikzpicture}
\end{titlepage}
